
\documentclass{article}
\usepackage{spconf,amsmath,graphicx}
\usepackage{cite}
\usepackage{amsmath,amssymb,amsfonts}
\usepackage{graphicx}
\usepackage{textcomp}
\usepackage{xcolor}
\usepackage{booktabs}
\usepackage{multirow}
\usepackage{microtype}
\usepackage{tikz}
\usepackage{makecell}
\usepackage{array}
\usetikzlibrary{positioning, arrows.meta, shapes.geometric}

%% ---- Notation macros ---------------------------------------
\newcommand{\zseq}{\mathbf{z}}          % embedding sequence (bold)
\newcommand{\zt}{z_t}                   % frame-level embedding
\newcommand{\zC}{z_{[\mathcal{C}]}}    % content token
\newcommand{\ztau}{z_{[\tau]}}          % compressed temporal token
\newcommand{\htau}{h_\tau}              % pre-bottleneck temporal token
\newcommand{\muspf}{\hat{S}}            % MSPF vector
\newcommand{\R}{\mathbb{R}}
\newcommand{\enc}{\mathcal{E}}          % audio encoder
\newcommand{\agg}{\mathcal{A}}          % temporal aggregator
\newcommand{\pred}{\mathcal{P}}         % trajectory predictor

%% ---- Additional notation macros -----------------------------
\newcommand{\dec}{\mathcal{D}}          % temporal decoder
\newcommand{\mask}{\mathcal{M}}         % mask set
\newcommand{\vseq}{\mathbf{v}}          % velocity sequence
\newcommand{\rseq}{\mathbf{r}}          % residual sequence
\newcommand{\zhat}{\hat{\mathbf{z}}}
\newcommand{\vhat}{\hat{\mathbf{v}}}
\newcommand{\rhat}{\hat{\mathbf{r}}}
\newcommand{\qseq}{\mathbf{q}}          % decoder query sequence
\newcommand{\za}{z_a}                   % anchor token
\newcommand{\pe}{\mathrm{PE}}
\newcommand{\adain}{\mathrm{AdaLN}}
\newcommand{\KL}{D_{\mathrm{KL}}}
\newcommand{\MI}{\hat{I}}
\newcommand{\dist}{\mathrm{dist}}

%% ---- Result placeholder ------------------------------------
\newcommand{\res}{\texttt{--}}

%% ---- Figure placeholder ------------------------------------
\usepackage{tikz}
\usetikzlibrary{positioning, arrows.meta, shapes.geometric}

% ---- Figure placeholder box --------------------------------
\newcommand{\figbox}[3]{%
  \begin{center}
    \begin{tikzpicture}
      \draw[dashed, gray!55, thick, rounded corners=3pt]
      (0,0) rectangle (#2,#3);
      \node[
        gray!55,
        font=\small\itshape,
        align=center,
        text width=\dimexpr#2cm-0.6cm\relax
      ] at (#2/2, #3/2) {#1};
    \end{tikzpicture}
  \end{center}
}

%% ============================================================
\begin{document}

\title{progAE : full-track music representation learning with latent trajectory compression}

\name{Anonymous Authors}
\address{Anonymous Affiliations\\
\texttt{anonymous@anon.org}}

\maketitle

\ninept

%% ============================================================

\begin{abstract}
  Music representation models operate on short audio excerpts
  (3--30\,s), yet songs span several minutes with rich temporal
  structure---evolving mood arcs, harmonic progressions, structural
  builds---that average pooling of chunk embeddings discards.  We
  present a self-supervised framework for learning a compact
  \emph{temporal token} for full-length recordings by compressing a
  sequence of frozen pre-trained audio embeddings through a
  variational bottleneck.  The central hypothesis is that temporal
  evolution can be represented by a single low-dimensional embedding
  that is complementary to, rather than redundant with, the standard
  mean-pooled content token.  To evaluate this hypothesis without
  requiring dense musical annotations, we introduce the \emph{Music
  Semantic Progress Function} (MSPF), a closed-form descriptor of the
  cumulative rate of semantic change across a recording.  We evaluate
  whether the temporal token reconstructs embedding trajectories,
  responds to order-changing perturbations, exposes structure-related
  information to lightweight probes, and remains distinct from the
  content token.  These experiments establish the temporal token as a
  compact representation for analysis, retrieval, and future
  generative conditioning.
\end{abstract}

%% ============================================================
\section{Introduction}
\label{sec:intro}

%% TODO — write in full after confirming experimental narrative.
%% Outline:
%%
%% Para 1 — The long-context gap
%%   Music models encode 3–30 s windows; songs are 3–5 minutes.
%%   Standard practice: slide encoder over track, average-pool.
%%   Average pooling is order-invariant — it cannot distinguish a
%%   song that builds to a drop from one that decays from the start.
%%
%% Para 2 — Why this matters in practice
%%   DJ tools, sync licensing, editorial workflows all require
%%   temporal event queries ("find tracks that build after 90 s").
%%   Structural segmentation, time-varying attribute tagging, and
%%   temporal retrieval are all bottlenecked by average pooling.
%%
%% Para 3 — Prior work and its limits
%%   Dyno [cite] introduced the aggregator/predictor paradigm for
%%   temporal tokens; demonstrated improvements on structural
%%   segmentation and temporal tagging.  Key limitation: no
%%   principled compression, no annotation-free evaluation metric
%%   (relied on SSMs, requiring reference reconstructions).
%%
%% Para 4 — The compression insight from video
%%   Stracke et al. [cite] learn motion embeddings in video using
%%   64× temporal compression and show quality improves with
%%   compression depth.  We test whether this transfers to music.
%%
%% Para 5 — The MSPF insight from video
%%   Metzer et al. [cite] introduce the Semantic Progress Function
%%   (SPF), a closed-form cumulative semantic distance that provides
%%   a model-agnostic diagnostic for temporal evolution in video.
%%   We adapt this to music as the MSPF.
%%
%% Para 6 — Our contributions (numbered)
%%   (1) Variational temporal bottleneck with compression ablation.
%%   (2) MSPF: the first annotation-free metric for temporal music
%%       representations.
%%   (3) Receptive-field analysis linking encoder choice to
%%       content-temporal entanglement.

\noindent\textcolor{gray!80}{\itshape[Introduction: to be written.]}

%% ============================================================
\section{Related Work}
\label{sec:related}

%% TODO — write in full.
%% Subsections:
%%
%% 2.1  Long-context music representation
%%      — Average pooling, hierarchical encoding [cite],
%%        attention aggregators [cite], Perceiver IO [cite].
%%      — Dyno [cite]: aggregator+predictor paradigm, SSM loss.
%%
%% 2.2  Variational and bottleneck representations
%%      — Beta-VAE [Higgins 2017, cite]: KL regularisation for
%%        disentangled latents.
%%      — Information bottleneck [Tishby, cite].
%%      — VQ-VAE for audio [cite].
%%
%% 2.3  Temporal representations in video
%%      — Motion embeddings [Stracke 2026, cite]: 64× compression
%%        improves quality; flow-matching generation in motion space.
%%      — Trajectory evaluation [Baumann 2026, cite]: Myriad /
%%        OWM benchmark; distributional accuracy + diversity.
%%      — Semantic Progress Function [Metzer 2026, cite]: cumulative
%%        semantic distance; retiming for linearisation.
%%
%% 2.4  Cover song identification
%%      — CQT-based [cite], learning-based [cite].
%%      — Used here as a controlled proxy for structure-sensitive,
%%        content-invariant temporal retrieval.
%%
%% 2.5  Music structure analysis
%%      — SALAMI [cite], boundary detection pipelines.
%%      — Novelty curves, self-similarity matrices as structure proxies.

\noindent\textcolor{gray!80}{\itshape[Related Work: to be written.]}

%% ============================================================
%% ============================================================
\section{Method}
\label{sec:method}

\subsection{Problem Formulation}
\label{ssec:prob}

Let $\enc$ be a frozen pre-trained audio representation model with
receptive field $T_s$\,s.  Given a full-length track of duration
$T_a$, we apply $\enc$ with stride $\Delta$ and obtain a sequence of
local music embeddings
%
\begin{equation}
  \zseq = (z_1, z_2, \ldots, z_T),
  \qquad
  z_t \in \R^D,
  \qquad
  T \approx \frac{T_a - T_s}{\Delta}.
  \label{eq:seq}
\end{equation}
%
Each $z_t$ summarizes the local musical content around time $t$,
but the full sequence $\zseq$ contains additional information about
how the track evolves: structural repetitions, builds, drops,
sectional boundaries, and long-range semantic arcs.  The standard
track-level representation is the temporal mean
%
\begin{equation}
  \zC = \frac{1}{T}\sum_{t=1}^{T} z_t,
  \qquad
  \zC \in \R^D,
  \label{eq:content}
\end{equation}
%
which aggregates local content but is order-invariant.  Any
permutation of $\zseq$ produces the same $\zC$, so average pooling
cannot distinguish, for example, a track that gradually builds into
a chorus from one with similar semantics that begins with high intensity and decays.

Our goal is to learn a compact temporal token
%
\begin{equation}
  \ztau \in \R^d,
  \qquad d \ll TD,
\end{equation}
%
that captures the temporal evolution of the embedding trajectory
while remaining complementary to the content token $\zC$.  We train
a self-supervised autoencoder over frozen embedding trajectories:
an aggregator compresses $\zseq$ into $\ztau$, and a conditional
Transformer decoder reconstructs the embedding trajectory through
residual prediction from $\ztau$, positional decoder inputs, and the
content anchor $\zC$.  Velocity prediction is retained as an
alternative training objective.

\subsection{Content Anchor and Reconstruction Targets}
\label{ssec:targets}

To prevent the temporal token from having to encode absolute
track-level content, the decoder is provided with an anchor
embedding $\za$.  We consider two anchors:
%
\begin{equation}
  \za = z_1
  \quad \text{(start-frame anchor)},
  \qquad
  \za = \zC
  \quad \text{(content anchor)}.
  \label{eq:anchor}
\end{equation}
%
The anchor defines the content coordinate system in which the
temporal token describes a trajectory.  Unless otherwise stated, we
use the mean content anchor $\za=\zC$.

We train the model to reconstruct either residuals or velocities.
The residual target is
%
\begin{equation}
  r_t = z_t - \za,
  \qquad t=1,\ldots,T,
  \label{eq:residual}
\end{equation}
%
while the velocity target is
%
\begin{equation}
  v_t = z_{t+1} - z_t,
  \qquad t=1,\ldots,T-1.
  \label{eq:velocity}
\end{equation}
%
Residual prediction directly reconstructs the trajectory relative
to the anchor.  Velocity prediction instead reconstructs local
semantic displacement and recovers the sequence by integration:
%
\begin{equation}
  \hat{z}_1 = \za,
  \qquad
  \hat{z}_{t+1} = \hat{z}_{t} + \hat{v}_{t}.
  \label{eq:vel_integration}
\end{equation}
%

\subsection{Temporal Aggregator}
\label{ssec:aggregator}

The temporal aggregator $\agg$ maps the full sequence of frozen
embeddings to a pre-bottleneck temporal summary.  By default, we
prepend a learnable CLS token $[\mathrm{CLS}] \in \R^D$ to the sequence and
process the result with a lightweight Transformer encoder:
%
\begin{equation}
  [z_{[\tau]}, h_1, \ldots, h_T]
  =
  \agg\!\left([\mathrm{CLS}, z_1, \ldots, z_T]\right),
  \label{eq:agg}
\end{equation}
%
The aggregator uses Pre-LayerNorm Transformer blocks, multi-head self-attention,
GELU/SwiGLU feed-forward layers, and RoPE positional encoding.

We optionally adopt a variational autoencoder formulation, where
the temporal token is parametrized as a Gaussian posterior
$\mathcal{N}(\mu_\tau,\operatorname{diag}(\sigma_\tau^2))$.  The
token is sampled using the reparameterisation trick:
%
\begin{equation}
  \ztau = \mu_\tau + \sigma_\tau \odot \varepsilon,
  \qquad
  \varepsilon \sim \mathcal{N}(0,I_d).
  \label{eq:reparam}
\end{equation}
%
At inference time, unless otherwise stated, we use the posterior
mean $\mu_\tau$ as the deterministic temporal representation for
retrieval, probing, and perturbation evaluation.

This aggressive compression is motivated by recent work on
long-term motion embeddings, where increasing temporal compression
improves latent semantic structure and downstream generation
quality~\cite{stracke2026learning}.  In our setting, we test
whether a similar principle holds for music: strong compression
should discard local content and noise while preserving the
large-scale temporal arc.

\subsection{RoPE Mask-Token Transformer Decoder}
\label{ssec:decoder}

The decoder reconstructs a full residual or velocity trajectory
from a sequence of positional decoder inputs.  We initialize a
learned mask token $m \in \R^{D_{\dec}}$ at each target timestep
and add rotary temporal position information:
%
\begin{equation}
  q_t^{(0)} = m + \pe_{\mathrm{RoPE}}(t),
  \qquad
  t=1,\ldots,T',
  \label{eq:decoder_time_inputs}
\end{equation}
%
where $T'=T$ for residual or frame prediction, and $T'=T-1$ for
velocity prediction.  These positional inputs specify \emph{where}
in the track the decoder should reconstruct.  The temporal token
$\ztau$ specifies \emph{which} temporal trajectory should unfold.

The temporal token is injected into each decoder block using
adaptive LayerNorm conditioning.  For hidden state $x$ in a decoder
block,
%
\begin{equation}
  \adain(x;\ztau)
  =
  \gamma(\ztau) \odot \mathrm{LN}(x)
  +
  \beta(\ztau),
  \label{eq:adaln}
\end{equation}
%
where $\gamma(\cdot)$ and $\beta(\cdot)$ are learned projections
from $\ztau$.  In the default AdaLN-Zero parameterisation, the
modulation layers are initialized such that the decoder begins near
an identity or weakly conditioned function, stabilising early
training.

The decoder is therefore a conditional temporal generator:
%
\begin{equation}
  \hat{\vseq}
  =
  \dec_\theta
  \left(
    \{m+\pe_{\mathrm{RoPE}}(t)\}_{t=1}^{T-1};
    \ztau,\za
  \right),
  \label{eq:decoder_velocity}
\end{equation}
%
for velocity prediction, or
%
\begin{equation}
  \hat{\rseq}
  =
  \dec_\theta
  \left(
    \{m+\pe_{\mathrm{RoPE}}(t)\}_{t=1}^{T};
    \ztau,\za
  \right),
  \label{eq:decoder_residual}
\end{equation}
%
for residual prediction.  The anchor $\za$ is supplied either by
concatenating a projected anchor embedding to each positional input, by
using a separate AdaLN modulation, or by cross-attending to an
anchor token; we use AdaLN by default and ablate other approaches.

This design differs from a standard sequence autoencoder in two
ways.  First, the decoder never receives the original input tokens;
it starts from position-aware mask tokens, forcing the trajectory to
be generated from the compressed temporal token.  Second, the
decoder reconstructs relative motion in a content coordinate system
defined by $\za$, encouraging $\ztau$ to encode temporal evolution
rather than static content.

\subsection{Encoder-Side Corruption and Span Masking}
\label{ssec:spanmask}

Decoder mask tokens alone do not prevent the aggregator from seeing
the full input sequence.  To further encourage global temporal
abstraction, we optionally corrupt the encoder input with temporal
masking:
%
\begin{equation}
  \tilde{\zseq}
  =
  \mathrm{Mask}(\zseq;\mathcal{M}),
  \qquad
  \ztau = \agg(\tilde{\zseq}),
  \label{eq:encoder_masking}
\end{equation}
%
where $\mathcal{M}$ is a set of masked timesteps.  We consider two
masking schemes.

\textbf{Random timestep masking} masks isolated timesteps sampled
uniformly across the track.  This tests robustness to missing
local observations but may be too easy because neighbouring music
embeddings are highly correlated.

\textbf{Contiguous span masking} masks temporally contiguous spans
of length $L_{\mathrm{span}}$.  This is the default corruption
strategy for the masked-input ablation because it forces the
aggregator to infer section-level evolution from incomplete
context.  We sweep
%
\begin{equation}
  p_{\mathrm{mask}} \in \{0, 0.25, 0.5, 0.75\},
  \qquad
  L_{\mathrm{span}} \in \{5,15,30,60\}\,\mathrm{s}.
  \label{eq:mask_sweep}
\end{equation}
%
The decoder always reconstructs the full residual or velocity
trajectory, including masked regions.

\subsection{Training Objective}
\label{ssec:train}

The default training objective combines residual reconstruction and
KL regularisation:
%
\begin{align}
  \mathcal{L}
  &=
  \lambda_r
  \underbrace{
    \frac{1}{T}
    \sum_{t=1}^{T}
    \left\|
    \hat{r}_t - (z_t-\zC)
    \right\|_2^2
  }_{\mathcal{L}_{\mathrm{res}}}
  \nonumber\\
  &\quad
  + \beta
  \underbrace{
    \KL\!\left(
      \mathcal{N}(\mu_\tau,\sigma_\tau^2)
      \,\|\, \mathcal{N}(0,I)
    \right)
  }_{\mathcal{L}_{\mathrm{KL}}}.
  \label{eq:default_loss}
\end{align}
%
For the velocity-prediction ablation, we replace
$\mathcal{L}_{\mathrm{res}}$ with
%
\begin{equation}
  \mathcal{L}_{\mathrm{vel}}
  =
  \frac{1}{T-1}
  \sum_{t=1}^{T-1}
  \left\|
  \hat{v}_t - (z_{t+1}-z_t)
  \right\|_2^2.
  \label{eq:vel_loss}
\end{equation}
%
For direct frame prediction, we use
%
\begin{equation}
  \mathcal{L}_{\mathrm{frame}}
  =
  \frac{1}{T}
  \sum_{t=1}^{T}
  \left\|
  \hat{z}_t - z_t
  \right\|_2^2.
  \label{eq:frame_loss}
\end{equation}
%

% For masked-input training, the reconstruction loss is computed over
% all timesteps, but we additionally report masked-region error:
% %
% \begin{equation}
%   \mathcal{L}_{\mathrm{mask}}
%   =
%   \frac{1}{|\mathcal{M}|}
%   \sum_{t\in\mathcal{M}}
%   \left\|
%     \hat{z}_t - z_t
%   \right\|_2^2.
%   \label{eq:mask_loss}
% \end{equation}
% %
% This term is used for evaluation and may optionally be added as a
% separate weighted training loss.

% Finally, for disentanglement ablations, we include an optional
% mutual-information penalty between temporal and content tokens:
% %
% \begin{equation}
%   \mathcal{L}_{\mathrm{dis}}
%   =
%   \MI(\ztau;\zC),
%   \label{eq:mi_loss}
% \end{equation}
% %
% estimated using MINE~\cite{belghazi2018mine} or a $k$-NN mutual
% information estimator~\cite{kraskov2004estimating}.  The complete
% optional objective is
%
% \begin{equation}
%   \mathcal{L}
%   =
%   \lambda_v \mathcal{L}_{\mathrm{vel}}
%   +
%   \lambda_r \mathcal{L}_{\mathrm{res}}
%   +
%   \lambda_m \mathcal{L}_{\mathrm{mask}}
%   +
%   \beta \mathcal{L}_{\mathrm{KL}}
%   +
%   \lambda_{\mathrm{dis}}\mathcal{L}_{\mathrm{dis}}.
%   \label{eq:full_loss}
% \end{equation}
%
The default model uses residual prediction with
$\lambda_r=1$, $\lambda_v=\lambda_m=\lambda_{\mathrm{dis}}=0$,
and $\beta=0.01$.

\subsection{Training Details}
\label{ssec:training_details}

Unless otherwise stated, models are trained on MTG-Jamendo
\cite{bogdanov2019mtg} using frozen MuQ embeddings
\cite{muq2025} extracted at 1\,Hz.  We use AdamW with learning rate
$10^{-4}$, cosine decay, 5\,000 warmup steps, batch size 64, and a
maximum of 200\,000 optimisation steps.  We retain the checkpoint
with lowest full-validation loss and stop after 10 validation checks
without improvement.  The default aggregator is a 6-layer
Transformer encoder with width 768, 12 attention heads, and
feed-forward dimension 3072.  The decoder has 4 layers, width 512,
8 attention heads, and feed-forward dimension 2048, with AdaLN
conditioning from $\ztau$.  Unless otherwise specified, $d=32$,
$\beta=0.01$, $\Delta=1$\,s, residual prediction, and the content
anchor $\za=\zC$ are used.

\begin{figure*}[t]
  \centering
  \includegraphics[width=0.8\textwidth]{figs/Schema.png}
  \caption{
    Overview of the proposed approach for learning dynamic information in a compressed representation. A pretrained audio encoder $\enc$ produces an embedding sequence $\zseq$; the temporal aggregator $\agg$ compresses this into a temporal token $\ztau$, while the content token $\zC$ is computed as a mean over time. The predictor $\pred$ reconstructs the original sequence $\hat{\zseq}$ from learned positional decoder inputs, conditioned on both $\zC$ and $\ztau$. The red arrow denotes the training objective between $\hat{\zseq}$ and $\zseq$.
  }

  \label{fig:method_overview}
\end{figure*}

%% %% ============================================================
%% ============================================================
\section{Experiments}
\label{sec:experiments}

We evaluate whether a compact temporal token captures the semantic
evolution of full-length music beyond order-invariant content
pooling.

Unless otherwise stated, all models are trained on MTG-Jamendo and
evaluated on held-out tracks.  The main experiments use frozen MuQ
embeddings extracted at 1\,Hz.  Retrieval and distance-based
experiments use cosine distance for the content token $\zC$ and
Euclidean distance for the temporal token $\ztau$ after
standardisation.

We use 51,701 tracks for training, 2,000 for validation, and 2,000
for testing.  Splits are disjoint by artist and album, use all 55,701
available extracted tracks, and are fixed across paper experiments.

The evaluation contract is fixed across the suite.  Reconstruction,
MSPF, perturbation, retrieval, and compression results use the same
held-out split and are reported at the track level.  Gain, pitch-shift,
and time-stretch perturbations are generated from the evaluation audio
before feature extraction.  Chunk shuffle, reversal, and section
deletion are applied directly to the frozen embedding sequence.  Unless
otherwise stated, MSPF metrics are
computed in MuQ space; cross-encoder MSPF checks recompute the same
metrics with MERT, USAD, and Music2Latent features.  Diffusion
experiments draw $N$ samples for each conditioning pair and report
both best-sample and distributional statistics so that stochastic
generation is not evaluated from a single cherry-picked trajectory.

\subsection{Datasets and Evaluation Settings}
\label{ssec:exp_setup}

We use MTG-Jamendo as the primary training corpus because it
contains full-length music recordings at scale.  Held-out
MTG-Jamendo tracks are used for trajectory reconstruction,
MSPF-based evaluation, perturbation sensitivity, and ablations.
For supervised validation of musical structure, we use datasets
with functional segment annotations, such as SALAMI and Harmonix.
These annotations provide section boundaries and function labels
such as intro, verse, chorus, bridge, instrumental, outro, and
silence.

For structure probing, we follow the standard formulation of music
structure analysis as two frame-level tasks: boundary detection and
function prediction.  Boundary detection identifies section changes,
while function prediction assigns a structural role to each frame.
Following prior work, we report boundary F-scores at 0.5\,s and
3\,s tolerance, denoted HR.5F and HR3F, and frame-level function
prediction metrics, including pairwise F-score (PWF) and accuracy
(ACC).

Figure~\ref{fig:mspf_qualitative} provides the qualitative
MSPF validation used throughout the experimental section.  It
visualises MSPF curves alongside musical structure annotations and
controlled temporal variants.

\begin{figure}[t]
  \centering
  \figbox{Qualitative MSPF validation.\\
    Placeholder: section labels, chord or structure annotations,\\
    and MSPF curves for original, cover, shuffle, reverse,\\
  and section-deletion variants.}{7.8}{3.0}
  \caption{Qualitative validation of the Music Semantic Progress
    Function (MSPF).  The figure illustrates how MSPF evolves over
    time for an original track and controlled temporal variants.  The
    purpose is to verify that MSPF follows large-scale musical change
  while reacting to transformations that alter temporal order.}
  \label{fig:mspf_qualitative}
\end{figure}

\subsection{Trajectory Reconstruction}
\label{ssec:exp_reconstruction}

The first experiment tests whether the compressed temporal token
preserves enough information to reconstruct the trajectory of frozen
music embeddings.  Given an input sequence $\zseq$, the model
encodes it into $\ztau$ and reconstructs the trajectory
$\hat{\zseq}$ using the decoder conditioned on $\ztau$ and an
anchor $z_a$.  This evaluates whether a single low-dimensional token
can represent the temporal evolution that is discarded by
order-invariant pooling.

We compare temporal-token models against content-only baselines.
The content-only decoder receives the same anchor information but no
temporal token, and therefore measures how much of the trajectory
can be reconstructed from static content alone.  We additionally
include a model trained to reconstruct from content-only
conditioning, which serves as a collapse diagnostic: if temporal
information is unavailable, the decoder should converge toward an
average or content-dominated trajectory.

We report two metrics.  First, per-frame reconstruction error
measures fidelity to the original embedding sequence:
%
\begin{equation}
  \mathcal{E}_{\mathrm{rec}}
  =
  \frac{1}{T}
  \sum_{t=1}^{T}
  \|\hat{z}_t-z_t\|_2^2.
  \label{eq:exp_recon_mse}
\end{equation}
%
Second, MSPF $R^2$ measures how much variance in semantic progress
is preserved by the reconstructed trajectory:
%
\begin{equation}
  R^2_{\mathrm{MSPF}}
  =
  1 -
  \frac{
    \sum_t (\muspf_t-\hat{\muspf}_t)^2
  }{
    \sum_t (\muspf_t-\bar{\muspf})^2
  }.
  \label{eq:exp_mspf_r2}
\end{equation}
%
The combination of these metrics separates local reconstruction
quality from preservation of the global semantic trajectory.
Table~\ref{tab:reconstruction} reports trajectory reconstruction
for content-only baselines and temporal-token models with different
bottleneck dimensions.

\begin{table}[t]
  \centering
  \setlength{\tabcolsep}{5.0pt}
  \caption{Trajectory reconstruction on held-out tracks.  Lower
    reconstruction error and higher MSPF $R^2$ indicate better
    preservation of the embedding trajectory and its semantic-progress
  profile.}
  \label{tab:reconstruction}
  \begin{tabular}{lcc}
    \toprule
    \textbf{Model}
    & $\mathcal{E}_{\mathrm{rec}}\downarrow$
    & MSPF $R^2\uparrow$ \\
    \midrule
    Content mean $\zC$ & \res & \res \\
    Content-only trained decoder & \res & \res \\
    \midrule
    $\ztau$, $d=1$ & \res & \res \\
    $\ztau$, $d=16$ & \res & \res \\
    $\ztau$, $d=32$ & \res & \res \\
    $\ztau$, $d=64$ & \res & \res \\
    \bottomrule
  \end{tabular}
\end{table}

\subsection{Annotation-Free Validation of MSPF}
\label{ssec:exp_mspf}

The MSPF is intended to provide an annotation-free descriptor of
temporal semantic progress.  Before using it as a reference signal,
we validate its behaviour under controlled transformations.  A
useful temporal-progress descriptor should remain similar when the
large-scale musical arc is preserved, but should change when the
ordering or structure of the track is disrupted.

We compare the MSPF of an original track to MSPFs computed from
several variants: cover versions, random chunk shuffles, reversals,
and section deletions.  Covers are used as qualitative or
case-study examples because they may preserve large-scale temporal
form while varying timbre, production, tempo, or key.  Shuffling,
reversal, and section deletion provide stronger controlled tests
because they explicitly alter temporal order or remove structural
material.

We report dynamic time warping distance between MSPF curves:
%
\begin{equation}
  d_{\mathrm{MSPF}}(i,j)
  =
  \mathrm{DTW}
  \left(
    \muspf(\zseq_i), \muspf(\zseq_j)
  \right).
  \label{eq:mspf_dtw}
\end{equation}
%
DTW is used because musically corresponding events may occur at
slightly different times, especially for covers or time-stretched
variants.  For transformations where absolute timing is important,
we additionally visualise the unwarped MSPF curves in
Figure~\ref{fig:mspf_qualitative}.  Table~\ref{tab:mspf_validation}
summarises the MSPF-DTW distances for transformations expected to
preserve or disrupt temporal progress.

\begin{table}[t]
  \centering
  \setlength{\tabcolsep}{5.2pt}
  \caption{MSPF validation under controlled transformations.  DTW
    distance measures disagreement between the MSPF of the original
  track and the transformed version.}
  \label{tab:mspf_validation}
  \begin{tabular}{lc}
    \toprule
    \textbf{Condition}
    & MSPF DTW$\downarrow$ \\
    \midrule
    Cover version & \res \\
    Pitch shift & \res \\
    Moderate time stretch & \res \\
    Gain change & \res \\
    \midrule
    & MSPF DTW$\uparrow$ \\
    \midrule
    Chunk shuffle & \res \\
    Reverse & \res \\
    Section deletion & \res \\
    % Hard splice & \res \\
    \bottomrule
  \end{tabular}
\end{table}

After validating MSPF as a temporal-progress descriptor, we test
whether the learned temporal-token geometry agrees with MSPF-based
similarity.  We compute pairwise temporal-token distances and
compare them to pairwise MSPF-DTW distances over held-out tracks.
Because retrieval and neighbourhood structure depend on ranking
rather than calibrated distances, we report Spearman correlation:
%
\begin{equation}
  \rho_{\tau,\mathrm{MSPF}}
  =
  \mathrm{corr}_{\mathrm{rank}}
  \left(
    \|\ztau^{(i)}-\ztau^{(j)}\|_2,
    d_{\mathrm{MSPF}}(i,j)
  \right).
  \label{eq:mspf_geometry}
\end{equation}
%
This experiment asks whether tracks that are close in temporal-token
space also have similar semantic-progress profiles.
Table~\ref{tab:mspf_geometry} reports this agreement for the
content token, MSPF features, the temporal token, and their
combination.

\begin{table}[t]
  \centering
  \setlength{\tabcolsep}{5.0pt}
  \caption{Agreement between representation geometry and MSPF-based
  temporal-progress distance.}
  \label{tab:mspf_geometry}
  \begin{tabular}{lc}
    \toprule
    \textbf{Representation}
    & Spearman $\rho_{\tau,\mathrm{MSPF}}\uparrow$ \\
    \midrule
    Random & \res \\
    Content token $\zC$ & \res \\
    MSPF feature baseline & \res \\
    Temporal token $\ztau$ & \res \\
    $[\zC,\ztau]$ & \res \\
    \bottomrule
  \end{tabular}
\end{table}

To reduce circularity, we repeat the MSPF validation with encoders
that are not used to train the default temporal token.  For each
held-out track we compute the default $\ztau$ from MuQ, then compute
MSPF curves and MSPF-DTW distances in an independent frozen encoder
space.  Agreement under this protocol indicates that the temporal
token is not merely matching idiosyncrasies of the MuQ feature
geometry.  Table~\ref{tab:mspf_cross_encoder} reports this
cross-encoder MSPF validation.

\begin{table}[t]
  \centering
  \setlength{\tabcolsep}{3.0pt}
  \caption{Cross-encoder MSPF validation.  The temporal token is
    trained from MuQ embeddings; MSPF metrics are recomputed using an
  independent scoring encoder.}
  \label{tab:mspf_cross_encoder}
  \begin{tabular}{lccc}
    \toprule
    \textbf{MSPF encoder}
    & MSPF $R^2\uparrow$
    & $\rho_{\tau,\mathrm{MSPF}}\uparrow$
    & $\rho_{\mathrm{sep}}\uparrow$ \\
    \midrule
    MuQ & \res & \res & \res \\
    MERT & \res & \res & \res \\
    USAD & \res & \res & \res \\
    Music2Latent & \res & \res & \res \\
    \bottomrule
  \end{tabular}
\end{table}

\subsection{Temporal Retrieval Case Studies}
\label{ssec:exp_temporal_retrieval}

We complement the quantitative MSPF-geometry test with qualitative
nearest-neighbour retrieval.  Given a query track, we retrieve nearest
neighbours using only $\zC$, only $\ztau$, and the concatenated
representation $[\zC,\ztau]$.  The expected behaviour is different
for each space: $\zC$ should favour similar timbre, genre, and
instrumentation, while $\ztau$ should favour similar long-range
evolution such as gradual build-up, drop-and-return form, or repeated
verse--chorus alternation.  The concatenated representation is used
as a practical retrieval mode when both static identity and temporal
shape matter.

Figure~\ref{fig:temporal_retrieval} visualises representative
queries by plotting the query MSPF curve beside the nearest-neighbour
curves returned by each representation.  Table~\ref{tab:temporal_retrieval}
summarises the same case studies with compact musical descriptors.

\begin{figure}[t]
  \centering
  \figbox{Temporal retrieval case studies.\\
    Placeholder: query MSPF and nearest-neighbour MSPFs from\\
  $\zC$, $\ztau$, and $[\zC,\ztau]$ retrieval.}{7.8}{3.0}
  \caption{Qualitative temporal retrieval.  The comparison checks
    whether $\ztau$ retrieves tracks with similar semantic-progress
    curves even when content-token retrieval favours more similar
  surface audio characteristics.}
  \label{fig:temporal_retrieval}
\end{figure}

\begin{table}[t]
  \centering
  \setlength{\tabcolsep}{3.0pt}
  \caption{Nearest-neighbour retrieval case studies.  Entries are
    qualitative descriptors of the top retrieved track under each
  representation.}
  \label{tab:temporal_retrieval}
  \begin{tabular}{llll}
    \toprule
    \textbf{Query} & $\zC$ NN & $\ztau$ NN & $[\zC,\ztau]$ NN \\
    \midrule
    Gradual build & \res & \res & \res \\
    Verse--chorus & \res & \res & \res \\
    Drop--return & \res & \res & \res \\
    Sparse outro & \res & \res & \res \\
    \bottomrule
  \end{tabular}
\end{table}

\subsection{Sensitivity to Temporal Perturbations}
\label{ssec:exp_perturbations}

A temporal representation should be stable under transformations
that preserve musical evolution, while responding to edits that
change temporal order or remove structural material.  We therefore
construct two families of perturbations.  The preserving family
includes pitch shift, moderate time stretch, gain change, and
optional generative edits such as instrument insertion when the
global arrangement is preserved.  The disrupting family includes
chunk shuffle, reversal, and section deletion, all applied in the
latent sequence domain so that they isolate temporal organisation
without introducing audio-encoder artefacts.

For each original track $x$ and perturbed track $x'$, we measure
normalised displacement in temporal-token space:
%
\begin{equation}
  \delta_\tau(x,x')
  =
  \frac{
    \|\ztau(x)-\ztau(x')\|_2
  }{
    \bar{d}_\tau
  },
  \label{eq:normalised_token_displacement}
\end{equation}
%
where $\bar{d}_\tau$ is the mean pairwise temporal-token distance
over the evaluation set.  We summarise the separation between
disrupting and preserving perturbations as
%
\begin{equation}
  \rho_{\mathrm{sep}}
  =
  \frac{
    \mathbb{E}_{x,x'\in\mathcal{D}_{\mathrm{disrupt}}}
    [\delta_\tau(x,x')]
  }{
    \mathbb{E}_{x,x'\in\mathcal{D}_{\mathrm{preserve}}}
    [\delta_\tau(x,x')]
  }.
  \label{eq:perturbation_separation}
\end{equation}
%
This ratio gives a compact measure of whether the representation
distinguishes temporal edits from nuisance changes.

For chunk shuffling, we additionally define a continuous severity
score based on the amount of temporal displacement induced by the
permutation $\pi$:
%
\begin{equation}
  s(\pi)
  =
  \frac{
    \sum_{t=1}^{T}|t-\pi(t)|
  }{
    \max_{\pi'}
    \sum_{t=1}^{T}|t-\pi'(t)|
  }.
  \label{eq:shuffle_severity}
\end{equation}
%
We then report the Spearman correlation between shuffle severity
and temporal-token displacement.  This tests whether the token
changes gradually with the magnitude of the temporal edit rather
than merely detecting that a transformation has occurred.
Figure~\ref{fig:perturbation_sensitivity} visualises this relation,
Table~\ref{tab:perturbations} reports displacement under preserving
and disrupting perturbations, and Table~\ref{tab:shuffle_severity}
reports the shuffle-severity correlation.

\begin{figure}[t]
  \centering
  \figbox{Perturbation sensitivity.\\
    Placeholder: temporal-token displacement versus shuffle severity,\\
  with preserving transformations shown as reference points.}{7.8}{3.0}
  \caption{Sensitivity of the temporal token to controlled
    perturbations.  The experiment measures whether displacement in
    temporal-token space increases with the severity of temporal
  reordering while remaining small for preserving transformations.}
  \label{fig:perturbation_sensitivity}
\end{figure}

\begin{table}[t]
  \centering
  \setlength{\tabcolsep}{3.5pt}
  \caption{Perturbation sensitivity.  Preserving transformations should
    produce small displacement, while disrupting transformations should
  produce larger displacement.}
  \label{tab:perturbations}
  \begin{tabular}{lccccc}
    \toprule
    \textbf{Rep.}
    & Pitch$\downarrow$
    & Stretch$\downarrow$
    & Gain$\downarrow$
    & Disrupt$\uparrow$
    & $\rho_{\mathrm{sep}}\uparrow$ \\
    \midrule
    $\zC$ & \res & \res & \res & \res & \res \\
    MSPF feat. & \res & \res & \res & \res & \res \\
    $\ztau$, $d=1$ & \res & \res & \res & \res & \res \\
    $\ztau$, $d=16$ & \res & \res & \res & \res & \res \\
    $\ztau$, $d=32$ & \res & \res & \res & \res & \res \\
    $\ztau$, $d=64$ & \res & \res & \res & \res & \res \\
    \bottomrule
  \end{tabular}
\end{table}

\begin{figure}[t]
  \centering

  \begin{minipage}{0.48\linewidth}
    \centering
    \figbox{Scatter: severity vs. displacement\\
      $d=1$, CR=\res\\
    $\rho_s=\res$}{3.55}{2.35}
    \vspace{-0.5em}
    {\scriptsize (a) Extreme compression}
  \end{minipage}
  \hfill
  \begin{minipage}{0.48\linewidth}
    \centering
    \figbox{Scatter: severity vs. displacement\\
      $d=32$, CR=\res\\
    $\rho_s=\res$}{3.55}{2.35}
    \vspace{-0.5em}
    {\scriptsize (b) Default bottleneck}
  \end{minipage}

  \vspace{0.6em}

  \begin{minipage}{0.48\linewidth}
    \centering
    \figbox{Scatter: severity vs. displacement\\
      $d=16$, CR=\res\\
    $\rho_s=\res$}{3.55}{2.35}
    \vspace{-0.5em}
    {\scriptsize (c) Moderate compression}
  \end{minipage}
  \hfill
  \begin{minipage}{0.48\linewidth}
    \centering
    \figbox{Scatter: severity vs. displacement\\
      $d=64$, CR=\res\\
    $\rho_s=\res$}{3.55}{2.35}
    \vspace{-0.5em}
    {\scriptsize (d) Weak compression}
  \end{minipage}

  \caption{Relationship between temporal perturbation severity and
    temporal-token displacement for different compression ratios.  Each
    point corresponds to a chunk-shuffled track, with severity computed
    from Eq.~\eqref{eq:shuffle_severity} and displacement computed from
    Eq.~\eqref{eq:normalised_token_displacement}.  The Spearman
    correlations shown in each panel correspond to the values reported
  in Table~\ref{tab:shuffle_severity}.}
  \label{fig:shuffle_severity_scatter}
\end{figure}

\begin{table}[t]
  \centering
  \setlength{\tabcolsep}{5.0pt}
  \caption{Correlation between chunk-shuffle severity and
  temporal-token displacement.}
  \label{tab:shuffle_severity}
  \begin{tabular}{lc}
    \toprule
    \textbf{Representation}
    & Spearman $\rho(s,\delta_\tau)\uparrow$ \\
    \midrule
    $\zC$ & \res \\
    MSPF feat. & \res \\
    $\ztau$, $d=1$ & \res \\
    $\ztau$, $d=16$ & \res \\
    $\ztau$, $d=32$ & \res \\
    $\ztau$, $d=64$ & \res \\
    \bottomrule
  \end{tabular}
\end{table}

\subsection{Structure Probing}
\label{ssec:exp_structure}

We next test whether the temporal token contains information useful
for annotated music structure analysis.  Instead of retrieving
tracks with similar symbolic form, which can be sparse and
ill-defined for small annotated datasets, we evaluate boundary
detection and function prediction using lightweight probes.

Inspired by Toyama et al., we freeze all representations and jointly
predict boundaries and structural functions.  We replace their
framewise linear head with a deliberately small, one-layer Transformer
probe that observes the complete track.  Each probe uses the exact
encoder and feature rate used to train the evaluated checkpoint, with
no probe-time resampling.

The local baseline projects the complete native sequence $z_{1:T}$
into the probe.  Token-only probes receive no local features.  Instead,
they begin from a learned mask token repeated over the $T$ frame
positions, add positional encodings, and condition the Transformer
block through adaptive layer normalisation.  In the combined variants,
the native local sequence supplies the input while the selected global
tokens condition the same block through adaptive layer normalisation.
Tracks are padded only for batching: padded positions are masked from
self-attention, excluded from both losses, and removed before metric
computation.
We compare:
%
\begin{align}
  u \in \{&z_{1:T},\; \zC,\; \ztau,\; [\zC,\ztau],
  \nonumber\\
  &[z_{1:T},\zC],\; [z_{1:T},\ztau],\;
  [z_{1:T},\zC,\ztau]\}.
  \label{eq:probe_inputs}
\end{align}
%
This design asks whether a global token can reconstruct a framewise
structural labelling when given only position and the token itself.
Comparing $\ztau$ with $z_{1:T}$ measures how much of the full
sequence's structural information survives temporal compression;
comparing $\ztau$ with $\zC$ tests whether that information depends on
temporal order.

Boundary detection is evaluated with HR.5F and HR3F.  Function
prediction is evaluated with PWF and ACC:
%
\begin{equation}
  \mathcal{L}_{\mathrm{probe}}
  =
  \mathcal{L}_{\mathrm{BCE}}^{\mathrm{boundary}}
  +
  \mathcal{L}_{\mathrm{CE}}^{\mathrm{function}}.
  \label{eq:probe_loss}
\end{equation}
%
The probe is not intended to establish state-of-the-art music
structure analysis.  Its purpose is to compare how much structural
information is recoverable from the frozen sequence and its compressed
global tokens under the same low-capacity decoder.

\begin{table}[t]
  \centering
  \setlength{\tabcolsep}{3.0pt}
  \caption{Full-track attention probing for music structure analysis.  Boundary
    detection is evaluated using HR.5F and HR3F; function prediction is
  evaluated using PWF and ACC.}
  \label{tab:structure_probe}
  \begin{tabular}{lcccc}
    \toprule
    \textbf{Probe input}
    & HR.5F$\uparrow$
    & HR3F$\uparrow$
    & PWF$\uparrow$
    & ACC$\uparrow$ \\
    \midrule
    Full sequence $z_{1:T}$ & \res & \res & \res & \res \\
    Content token $\zC$ & \res & \res & \res & \res \\
    Temporal token $\ztau$ & \res & \res & \res & \res \\
    $[\zC,\ztau]$ & \res & \res & \res & \res \\
    \midrule
    $[z_{1:T},\zC]$ & \res & \res & \res & \res \\
    $[z_{1:T},\ztau]$ & \res & \res & \res & \res \\
    $[z_{1:T},\zC,\ztau]$ & \res & \res & \res & \res \\
    \bottomrule
  \end{tabular}
\end{table}

\subsection{Content--Temporal Disentanglement}
\label{ssec:exp_disentanglement}

We evaluate whether $\ztau$ captures temporal order information not
available to mean pooling.  The first test uses exact sequence
permutations.  Given an embedding sequence $\zseq$, we construct
permuted variants $\pi(\zseq)$ using random shuffle, reversal,
half-swap, circular shift, and local shuffle.  These transformations
preserve the multiset of frame embeddings, and therefore leave the
mean-pooled content token unchanged up to numerical precision:
%
\begin{equation}
  \Delta_C
  =
  \|\zC(\zseq)-\zC(\pi(\zseq))\|_2,
  \qquad
  \Delta_\tau
  =
  \|\ztau(\zseq)-\ztau(\pi(\zseq))\|_2.
  \label{eq:order_sensitivity}
\end{equation}
%
A representation of temporal evolution should produce larger
$\Delta_\tau$ for order-altering transformations, while $\Delta_C$
serves as an order-invariant reference.  Table~\ref{tab:order_sensitivity}
reports the content-token and temporal-token displacements under
exact sequence permutations.

\begin{table}[t]
  \centering
  \setlength{\tabcolsep}{5.0pt}
  \caption{Order sensitivity under exact sequence permutations.  Mean
    pooling is invariant to permutations, while a temporal token should
  respond to order changes.}
  \label{tab:order_sensitivity}
  \begin{tabular}{lcc}
    \toprule
    \textbf{Transform}
    & $\Delta_C\downarrow$
    & $\Delta_\tau\uparrow$ \\
    \midrule
    Random shuffle & \res & \res \\
    Reverse & \res & \res \\
    Half-swap & \res & \res \\
    Circular shift & \res & \res \\
    Local shuffle & \res & \res \\
    \bottomrule
  \end{tabular}
\end{table}

Finally, we quantify similarity between content and temporal
representations using representation-level dependence measures.  We
report linear CKA or CKNNA between $\ztau$ and $\zC$, together with
two references: a same-representation upper reference and a
shuffled-pair lower reference.  These values are used comparatively
rather than as absolute estimates of independence.
Table~\ref{tab:dependence} reports these dependence measures across
reference comparisons and temporal-token variants.

\begin{table}[t]
  \centering
  \setlength{\tabcolsep}{4.5pt}
  \caption{Representation dependence between temporal and content
    tokens.  Shuffled pairs provide a lower reference; identical or
  content-like representations provide an upper reference.}
  \label{tab:dependence}
  \begin{tabular}{lcc}
    \toprule
    \textbf{Comparison}
    & CKA$\downarrow$
    & CKNNA$\downarrow$ \\
    \midrule
    $\zC$ vs. $\zC$ & \res & \res \\
    $\ztau$ vs. shuffled $\zC$ & \res & \res \\
    $\ztau$ vs. $\zC$ & \res & \res \\
    $\ztau$ vs. $\zC$, no anchor & \res & \res \\
    $\ztau$ vs. $\zC$, VAE bottleneck & \res & \res \\
    \bottomrule
  \end{tabular}
\end{table}

\subsection{Generative Conditioning and Token Swapping}
\label{ssec:exp_diffusion_swapping}

The strongest test of complementarity is generative: if $\zC$ and
$\ztau$ separate content from temporal evolution, then a conditional
generator should follow the content source when $\zC$ is swapped and
the temporal source when $\ztau$ is swapped.  We therefore train a
small diffusion model in the frozen embedding-trajectory space.  The
model denoises a trajectory while conditioning on both global tokens:
%
\begin{equation}
  \tilde{\zseq}
  =
  G_\phi(\epsilon;\zC,\ztau),
  \qquad
  \epsilon \sim \mathcal{N}(0,I).
  \label{eq:diffusion_conditioning}
\end{equation}
%
At test time, we decode matched and swapped conditions from two
tracks $A$ and $B$:
%
\begin{align}
  &(\zC(A),\ztau(A)),\quad (\zC(B),\ztau(B)),
  \nonumber\\
  &(\zC(A),\ztau(B)),\quad (\zC(B),\ztau(A)).
  \label{eq:diffusion_swaps}
\end{align}
%
For each condition $c$ we draw $N$ generated trajectories
$\{\tilde{\zseq}^{(n)}_c\}_{n=1}^{N}$.  Content alignment is
measured by comparing the mean generated embedding
$\bar{\tilde{z}}^{(n)}_c$ to the requested content token in
$\zC$-space.  Temporal alignment is measured by comparing the MSPF
of each generated trajectory to the MSPF of the requested temporal
source; an additional temporal-token alignment is computed by
re-encoding $\tilde{\zseq}^{(n)}_c$ with $\agg$.  We report
best-of-$N$ reconstruction error, mean reconstruction error,
sample diversity, and minimum/mean MSPF error:
%
\begin{align}
  \mathcal{E}_{\mathrm{best}}(c)
  &=
  \min_n \mathcal{E}_{\mathrm{rec}}
  (\tilde{\zseq}^{(n)}_c,\zseq_{\mathrm{target}}),
  \nonumber\\
  \mathcal{E}_{\mathrm{div}}(c)
  &=
  \frac{2}{N(N-1)}
  \sum_{i<j}
  \dist(\tilde{\zseq}^{(i)}_c,\tilde{\zseq}^{(j)}_c).
  \label{eq:diffusion_distribution_metrics}
\end{align}
%
This experiment directly measures whether $\zC$ controls static
musical identity while $\ztau$ controls the evolution of the
generated track, and whether the conditional generator covers a
distribution of plausible trajectories rather than a single averaged
path.  Table~\ref{tab:diffusion_swapping} reports the full
distributional swap-alignment protocol.

\begin{table*}[t]
  \centering
  \setlength{\tabcolsep}{3.0pt}
  \caption{Diffusion token swapping.  A successful split should make
    generated content follow $\zC$ and generated temporal evolution
    follow $\ztau$, including under crossed conditions.  Best and mean
  errors are computed over $N$ samples per condition.}
  \label{tab:diffusion_swapping}
  \begin{tabular}{lccccccc}
    \toprule
    \textbf{Condition}
    & $C$ align.$\uparrow$
    & $\tau$ align.$\uparrow$
    & Best err.$\downarrow$
    & Mean err.$\downarrow$
    & Div.$\uparrow$
    & MSPF min$\downarrow$
    & MSPF mean$\downarrow$ \\
    \midrule
    $(\zC(A),\ztau(A))$ & \res & \res & \res & \res & \res & \res & \res \\
    $(\zC(B),\ztau(B))$ & \res & \res & \res & \res & \res & \res & \res \\
    $(\zC(A),\ztau(B))$ & \res & \res & \res & \res & \res & \res & \res \\
    $(\zC(B),\ztau(A))$ & \res & \res & \res & \res & \res & \res & \res \\
    \bottomrule
  \end{tabular}
\end{table*}

\subsection{Compression Summary}
\label{ssec:exp_compression_summary}

Compression is a central experimental axis, but the main paper uses
it only to choose and justify a single-token operating point.  We
sweep the temporal bottleneck dimension $d$ while keeping one
temporal token.  The compression ratio is
%
\begin{equation}
  \mathrm{CR}=\frac{TD}{d}.
  \label{eq:compression_ratio}
\end{equation}
%
Table~\ref{tab:compression} summarises the trade-off between
trajectory fidelity, temporal utility, and content leakage.  All
variants use the same 200\,000-step ceiling and validation-loss
early-stopping rule.  The training-objective, anchor/target, masking,
and multi-token capacity ablations are reported in the appendix
because they refine the design rather than establish the core claim
that a single $\ztau$ can encode track-level temporal evolution.

\begin{table*}[t]
  \centering
  \setlength{\tabcolsep}{3.0pt}
  \caption{Single-token compression sweep.  The default model uses
    $d=32$ as a stable middle-capacity operating point for the main
    experiments.}
  \label{tab:compression}
  \begin{tabular}{cccccc}
    \toprule
    $d$ & CR
    & $\mathcal{E}_{\mathrm{rec}}\downarrow$
    & MSPF $R^2\uparrow$
    & $\rho_{\mathrm{sep}}\uparrow$
    & CKA$(\ztau,\zC)\downarrow$ \\
    \midrule
    1 & \res & \res & \res & \res & \res \\
    4 & \res & \res & \res & \res & \res \\
    16 & \res & \res & \res & \res & \res \\
    32 & \res & \res & \res & \res & \res \\
    64 & \res & \res & \res & \res & \res \\
    256 & \res & \res & \res & \res & \res \\
    \bottomrule
  \end{tabular}
\end{table*}

\subsection{Encoder Receptive Field}
\label{ssec:exp_receptive_field}

The frozen encoder defines the local units from which the temporal
token is built.  We therefore evaluate whether encoder receptive
field and frame rate affect temporal-content disentanglement.  We
train the same temporal bottleneck over multiple frozen encoders
with different temporal resolutions and context lengths.  Shorter
receptive fields provide more local observations, while longer
receptive fields may already mix local content and temporal context
before the aggregator.

For each encoder, we report temporal utility and content leakage
using the same downstream metrics as above.  This experiment tests
whether the quality of $\ztau$ depends primarily on the bottleneck
architecture or also on the temporal granularity of the frozen
features.  Table~\ref{tab:encoder_rf} reports the effect of frozen
encoder receptive field and frame rate.

\begin{table}[t]
  \centering
  \setlength{\tabcolsep}{3.5pt}
  \caption{Effect of frozen encoder receptive field and frame rate.
  The same temporal-token model is trained on each encoder.}
  \label{tab:encoder_rf}
  \begin{tabular}{lcccc}
    \toprule
    \textbf{Encoder}
    & \textbf{RF}
    & MSPF $R^2\uparrow$
    & $\rho_{\mathrm{sep}}\uparrow$
    & CKA$(\ztau,\zC)\downarrow$ \\
    \midrule
    MuQ & $\sim$10\,s & \res & \res & \res \\
    MERT & $\sim$10\,s & \res & \res & \res \\
    USAD & $\sim$1\,s & \res & \res & \res \\
    Music2Latent & var. & \res & \res & \res \\
    MATPAC & $\sim$1\,s & \res & \res & \res \\
    \bottomrule
  \end{tabular}
\end{table}

%% ============================================================
\bibliographystyle{IEEEtran}
\bibliography{refs}

\clearpage

% \appendix

% \section{Deferred Experiment Rationale}
% \label{app:deferred_rationale}

% The main paper keeps experiments that directly support the core
% claim: a single compact temporal token $\ztau$ captures track-level
% evolution not contained in the content token $\zC$.  We therefore
% keep trajectory reconstruction, MSPF validation, perturbation
% sensitivity, shuffle-severity analysis, structure probing,
% content--temporal dependence, diffusion-based token swapping,
% encoder receptive-field analysis, and a single-token compression
% sweep in the main text.  The following experiments are deferred
% because they refine the design rather than establish the core
% ICASSP narrative: multi-token capacity sweeps, bottleneck-objective
% ablations, anchor/target ablations, and encoder-side masking.

% \section{Additional Details on the MSPF}
% \label{app:mspf_details}

% \section{Additional Design Ablations}
% \label{app:compression_ablations}

% \subsection{Extended Compression Sweep}
% \label{ssec:exp_ablations}

% Compression is a central experimental axis.  We sweep the temporal
% bottleneck dimension $d$ and, when applicable, the number of latent
% tokens $K$.  The compression ratio is
% %
% \begin{equation}
%   \mathrm{CR}
%   =
%   \frac{TD}{Kd}.
%   \label{eq:compression_ratio_extended}
% \end{equation}
% %
% This experiment tests whether temporal abstraction improves under a
% strong bottleneck, and whether increasing capacity primarily
% improves reconstruction or also improves semantic temporal metrics.

% We report reconstruction MSE, MSPF $R^2$, perturbation separation,
% structure probing performance, and dependence with the content
% token.  These metrics distinguish three behaviours: faithful
% trajectory reconstruction, useful temporal abstraction, and content
% leakage.  Table~\ref{tab:compression_extended} reports the compression sweep
% across bottleneck dimensions and numbers of latent tokens.

% \begin{table}[t]
%   \centering
%   \setlength{\tabcolsep}{3.0pt}
%   \caption{Compression sweep.  The table compares reconstruction,
%     semantic-progress preservation, perturbation sensitivity, structure
%   probing, and dependence with content.}
%   \label{tab:compression_extended}
%   \begin{tabular}{cccccc}
%     \toprule
%     $K$ & $d$ & CR
%     & MSPF $R^2\uparrow$
%     & $\rho_{\mathrm{sep}}\uparrow$
%     & CKA$(\ztau,\zC)\downarrow$ \\
%     \midrule
%     1 & 1 & \res & \res & \res & \res \\
%     1 & 4 & \res & \res & \res & \res \\
%     1 & 16 & \res & \res & \res & \res \\
%     1 & 64 & \res & \res & \res & \res \\
%     1 & 256 & \res & \res & \res & \res \\
%     2 & 4 & \res & \res & \res & \res \\
%     4 & 4 & \res & \res & \res & \res \\
%     8 & 4 & \res & \res & \res & \res \\
%     \bottomrule
%   \end{tabular}
% \end{table}

% We then isolate the effect of the bottleneck objective.  The
% deterministic autoencoder tests compression without stochastic
% regularisation.  Variational models with different $\beta$ values
% test whether KL regularisation improves abstraction or merely
% reduces capacity.  Additional noisy or dropout bottlenecks are used
% as simple regularisation baselines.  Table~\ref{tab:bottleneck_ablation}
% reports the bottleneck objective ablation.

% \begin{table}[t]
%   \centering
%   \setlength{\tabcolsep}{4.0pt}
%   \caption{Bottleneck objective ablation.  The goal is to separate the
%     effect of compression from the effect of stochastic or variational
%   regularisation.}
%   \label{tab:bottleneck_ablation}
%   \begin{tabular}{lccc}
%     \toprule
%     \textbf{Config.}
%     & MSPF $R^2\uparrow$
%     & $\rho_{\mathrm{sep}}\uparrow$
%     & HR3F$\uparrow$ \\
%     \midrule
%     AE, $\beta=0$ & \res & \res & \res \\
%     VAE, $\beta=.001$ & \res & \res & \res \\
%     VAE, $\beta=.01$ & \res & \res & \res \\
%     VAE, $\beta=.1$ & \res & \res & \res \\
%     Fixed noise & \res & \res & \res \\
%     Dropout & \res & \res & \res \\
%     \bottomrule
%   \end{tabular}
% \end{table}

% We also ablate reconstruction targets and anchors.  Direct frame
% prediction asks the token to reconstruct absolute embeddings.
% Residual prediction asks it to reconstruct deviations from an
% anchor.  Velocity prediction asks it to reconstruct local semantic
% motion.  Comparing these objectives tests which formulation most
% directly encourages the token to encode temporal change rather than
% static content.  Table~\ref{tab:target_anchor_ablation} reports the
% target and anchor ablation.

% \begin{table}[t]
%   \centering
%   \setlength{\tabcolsep}{3.6pt}
%   \caption{Prediction target and anchor ablation.}
%   \label{tab:target_anchor_ablation}
%   \begin{tabular}{llccc}
%     \toprule
%     \textbf{Target} & \textbf{Anchor}
%     & $\mathcal{E}_{\mathrm{rec}}\downarrow$
%     & MSPF $R^2\uparrow$
%     & CKA$(\ztau,\zC)\downarrow$ \\
%     \midrule
%     Frame & none & \res & \res & \res \\
%     Frame & $z_1$ & \res & \res & \res \\
%     Residual & $z_1$ & \res & \res & \res \\
%     Residual & $\zC$ & \res & \res & \res \\
%     Velocity & $z_1$ & \res & \res & \res \\
%     Velocity & $\zC$ & \res & \res & \res \\
%     \bottomrule
%   \end{tabular}
% \end{table}

% Finally, we test whether encoder-side masking improves temporal
% abstraction.  Random masking removes isolated timesteps, while span
% masking removes contiguous regions.  The decoder always reconstructs
% the full trajectory.  This ablation tests whether forcing the
% aggregator to infer missing regions encourages more global temporal
% representations.  Table~\ref{tab:masking_ablation} reports the
% masking ablation.

% \begin{table}[t]
%   \centering
%   \setlength{\tabcolsep}{3.4pt}
%   \caption{Encoder-side masking ablation.  Masked-region error measures
%   imputation quality; the remaining metrics measure temporal utility.}
%   \label{tab:masking_ablation}
%   \begin{tabular}{lccc}
%     \toprule
%     \textbf{Masking}
%     & Mask err.$\downarrow$
%     & MSPF $R^2\uparrow$
%     & $\rho_{\mathrm{sep}}\uparrow$ \\
%     \midrule
%     None & --- & \res & \res \\
%     Random, $p=.25$ & \res & \res & \res \\
%     Random, $p=.50$ & \res & \res & \res \\
%     Span, $p=.25$, 15\,s & \res & \res & \res \\
%     Span, $p=.50$, 15\,s & \res & \res & \res \\
%     Span, $p=.50$, 30\,s & \res & \res & \res \\
%     Span, $p=.50$, 60\,s & \res & \res & \res \\
%     \bottomrule
%   \end{tabular}
% \end{table}

\end{document}
